\section{A separate degree ten companion in 212 integers}
\label{sec:reciprocal22}
The reciprocal companion is not a quartic. It has exactly 212 unrestricted
integer coordinates, 151 residuals, exact total degree 10, 55820 collected
monomials and coefficient height 31104. Its largest residual has degree
five and coefficient height 108. The polynomial is the sum of the squares
of its complete literal residual array; its variable order and coefficients
are supplied separately from the unchanged principal quartic.

The incoming reciprocal-pooling theorem comes from source 168. Its old
factorial-based sign bootstrap is not imported into the present factorial-free
front. Instead we prove a cofinal restricted chart in which both bases are
nonzero before the power requests are interpreted; unit bases are allowed
by the reciprocal theorem and are eliminated only after source restoration. The original
$H=2+4(C+\sum A_i)$ remains unchanged.

\subsection{A cofinal radix chart with no early sign assumption}
Restrict the source by the literal substitution
\[
 q=3z+1,\qquad z\in\Z.
\]
At every integer input either $q\ge1$, $q+1\ge2$, or $q\le-2$,
$q+1\le-1$. Both bases are therefore nonzero and have the same sign,
before any semantic power interpretation. The possible values $q=1$
and $q+1=-1$ are not prematurely excluded. This is an integral chart
of a restricted subfamily, not a unimodular change on all old radix values.
We must prove that the restriction preserves completeness.

Choose the horizon $t=32+4a^2$ with $a\not\equiv0\pmod3$, for example
$a=3v+1$ arbitrarily large. Then $12\mid t$, and these horizons remain
cofinal. Put $K=32t^4$ and $F=t^K$, in that order; the power is not $K^t$.
For $p=2,3$ set $v_p(F)=v$ and $v_p(t!)=a_p$. Then
$v>\max_{1\le i\le t}v_p(i)$. Expanding the numerator of the binomial gives
\[
 t!(B-1)=\prod_{i=1}^t(F+i)-t!
 =\sum_{\varnothing\ne S\subseteq\{1,\ldots,t\}}
       F^{|S|}\prod_{i\notin S}i.
\]
Each summand has valuation
\[
 a_p+\sum_{i\in S}\bigl(v-v_p(i)\bigr)\ge a_p+1.
\]
Canceling the known $p$-power and the remaining unit part of $t!$ gives
$B\equiv1\pmod p$. This is an integer valuation proof, not reduction of
uncanceled rational factors modulo $p$. In particular $\gcd(B,6)=1$.

As in the main completeness proof, choose the finite path, then sufficiently
large $y=Fw$ with odd $w$ and $t!\mid y$. Every beta modulus
$1+(i+1)y$ is odd. Their prescribed CRT class therefore has arbitrarily
large even representatives; choose an even $x$ satisfying all retained
bounds. Now $N=B^x$ is coprime to six. Fix the eight ordinary representatives
in their prescribed classes modulo $N^2$. Because $\gcd(N^2,6)=1$, the
remaining class for $C$ can be chosen even and in either of the two residues
modulo three for which $3\nmid H$. Its valid values form a cofinal progression
of step $6N^2$. Choose $C$ above all other nominal indices and representatives.
This preserves all previous CRT congruences and $B\mid C+F+1$.

Only now define $H$ and $q=(BH)^x$. Here $H\equiv2\pmod4$, $3\nmid H$,
$B\equiv1\pmod3$, and $x$ is positive even. Consequently
$q\equiv0\pmod4$ and $q\equiv1\pmod3$, hence $q\equiv4\pmod{12}$.
In particular $q\equiv1\pmod3$, so the coordinate $z=(q-1)/3$
is integral. The stronger modulo-twelve conclusion is useful only for
constructing canonical witnesses, not imposed on arbitrary source zeros. Finally choose the actual square
frame, Pell center and even-anchor cap witnesses as before. The choice
order has no cyclic dependence:
\[
 t,K,F,\text{path},y,x,B,N,\text{other representatives},C,H,q,\text{frame},A.
\]
Evenness of $x$ is a canonical completeness choice, not an assumption on
arbitrary soundness witnesses.

\subsection{The reciprocal polynomial and its exact power theorem}
For $n_0=q,n_1=q+1$ let $I_j$ denote the two request groups, containing
three and ten requests. Retain the actual canonical Pell pairs $(X_i,Y_i)$,
outputs $m_i$, indices $k_i$, and wires $Z_i=AY_i$. Set
\begin{align*}
 T_j&=2An_j-n_j^2-1,&H_j&=\sum_{i\in I_j}Y_i(m_i-X_i),\\
 S_j&=\sum_{i\in I_j}\bigl((m_i-X_i+Z_i)^2-Y_i^2\bigr),
 &E_j&=S_j-2n_jH_j.
\end{align*}
Delete the four private old pool coordinates and four rows, introducing
one unrestricted integer $\zeta$ and the single residual
\begin{equation}\label{eq:reciprocal}
 E_0T_1+E_1T_0-\zeta T_0T_1=0.
\end{equation}
It has degree at most five. All other residuals remain quadratic.

\begin{theorem}
Under the retained canonical-Pell and square-frame hypotheses,
\eqref{eq:reciprocal} is solvable in $\zeta\in\Z$ if and only if all
thirteen outputs are the exact powers $m_i=n_j^{k_i}$. An index is
nonnegative whenever its base has absolute value at least two; unit bases
retain their canonical signed powers. No coprimality of $T_0,T_1$ is required.
\end{theorem}
\begin{proof}
Write $\delta_i=m_i-X_i+(A-n_j)Y_i$ for $i\in I_j$.
Exact expansion gives
\[
 \sum_{i\in I_j}\delta_i^2=E_j-T_j\sum_{i\in I_j}Y_i^2.
\]
Thus~\eqref{eq:reciprocal} says that $\sum_{j,i}\delta_i^2/T_j$ is an
integer. Let $b=\lfloor\sqrt{W-1}\rfloor\ge16$ and
$C_*=[q(q+1)]^b\ne0$. Every signed index satisfies $|k_i|\le b$.
The integer
\[
 d_i=C_*m_i-(C_*/n_j^b)n_j^{b+k_i}
\]
is defined without a negative exponent. Canonical Pell congruences imply
$C_*\delta_i\equiv d_i\pmod{T_j}$, because $n_j$ is a unit modulo its
own $T_j\equiv-1\pmod{n_j}$. Therefore
\[
 \sum_{j=0}^1\sum_{i\in I_j}\frac{d_i^2}{T_j}\in\Z.
\]
We only multiplied by $C_*^2$ and changed a numerator by a multiple of its
own denominator. Neither $C_*$ nor the other modulus was inverted.

The frame gives $|n_j|,|m_i|,|k_i|\le b$. Hence
\[
 |C_*|\le b^{2b},\quad |C_*/n_j^b|\le b^b,\quad
 |d_i|\le b^{2b+1}+b^{3b}\le2b^{3b}.
\]
With $r=13\le15$,
\[
 \sum_i d_i^2\le4rb^{6b}<W^{3b+1}\le W^W
       <\min(|T_0|,|T_1|).
\]
Indeed $4r<W$, $b^2<W$, $3b+1\le b^2+1\le W$, and
$|T_j|\ge2|A|-W>W^W$. Also $\operatorname{sgn}(T_j)=
\operatorname{sgn}(An_j)$, so the denominators have the same sign.
The displayed integral reciprocal sum has absolute value below one and
summands of a common sign. It is zero and every $d_i=0$. Dividing the
identity in $\Q$ by the nonzero $C_*$ gives $m_i=n_j^{k_i}$.
A negative index at a nonunit base would force an integer times a
positive power of its absolute value to equal one, impossible.

Conversely exact powers give $T_j\mid\delta_i$ for every request.
The square identity gives $T_j\mid E_j$, and
$\zeta=E_0/T_0+E_1/T_1$ is an integral solution of
\eqref{eq:reciprocal}.
\end{proof}
The common sign and strict bound are essential. With opposite denominators,
nonzero squares can cancel; without the bound a nonzero reciprocal sum can
be an integer. The proof uses neither phenomenon.

\subsection{Restoration and the consistency equivalence}
At a new zero, first the unchanged signed cap and orientation recover the
canonical pairs. The affine radix chart already gives two nonzero bases,
so the theorem applies without invoking an unauthenticated radix power.
Restore the old coordinates by
\[
 L_j=E_j/T_j,\qquad K_j=2H_j+(2A-n_j)L_j.
\]
They are integral on this zero set, and both old rows
$K_j-2H_j-(2A-n_j)L_j=0$ and $S_j-n_jK_j+L_j=0$ vanish.
Only now apply the accepted old source bootstrap. Its retained singleton
$BH$ request, nonzero even $H$, odd nonzero $B$ and endpoint argument
exclude $x=0$ and hence $q=1$. The recovered $q$ powers then give
$y>0$ and the large endpoint bound on $x$, excluding $q+1=-1$. Thus
both bases are nonunits at every actual zero, the representative signs
are authenticated in the old acyclic order, and a finite syntactic
refutation is decoded. The reciprocal step itself did not assume these
later conclusions.
Conversely the cofinal construction above supplies a source zero in the
restricted chart, and $\zeta=L_0+L_1$ supplies the new zero.
These are witness-existence maps, not a polynomial automorphism of the
whole ambient lattice.

The finite compiler checks that the four deleted coordinates are private
to the four protected rows, retains every singleton request and every
norm/orientation/cap equation, and substitutes $q=3z+1$ throughout the
entire retained source and every actual frame atom. Protected transported
operations and two integral radical projections are rechecked against
that source. The literal degree-five row and its complete squared
expansion determine the stated exact degree 10 and coefficients. This
separate composition proves the same consistency equivalence on
$\Z^{212}$ without changing the principal quartic or claiming a global
variable minimum.
